\chapter{راهکار پیشنهادی}\label{Chap:Chap3}
\minitoc
\setcounter{footnote}{0}
\pagestyle{plain}

همانطور که در فصل قبل مطرح شد، مسئله‌ی دسته‌بندی بزرگ مقیاس، نمونه محدود زوایای مختلفی دارد که کار‌های پیشین، هر کدام بر روی قسمت خاصی از حل چالش این مسئله متمرکز بوده‌اند. در نتیجه استفاده‌ی مستقیم از هر یک دارای نقاط ضعفی در این مسئله است. در این فصل، با الهام از فواید و نقصان‌های هر کدام از این روش‌ها، یک راه‌حل جامع برای حل مسئله‌ی دسته‌بندی بزرگ مقیاس، نمونه محدود ارائه می‌دهیم. در ابتدا به نمادگذاری مسئله و بیان پیش‌نیاز‌ها می‌پردازیم.

\section{نمادگذاری مسئله و پیش‌نیازها}

در مسئله‌ی دسته‌بندی بزرگ مقیاس، نمونه محدود به دو مجموعه 
\trans{دادگان}{Dataset}
بزرگ‌مقیاس 
$\mathcal{D}_{source}$
و 
$\mathcal{D}_{target}$
دسترسی داریم. از این دو مجموعه به ترتیب با نام‌های دادگان مبدا و مقصد یا دادگان فاز 
\trans{متاآموزش}{Meta-Train}
 و
 \trans{متاآزمون}{Meta-Test}
  یاد میکنیم. هر یک از این مجموعه دادگان که به صورت
  $\mathcal{D}=\{(x^i, y_{super}^i, y_{base}^i)\}_{i=1}^{|\mathcal{D}|}$
  هستند، متشکل از سه مجموعه می‌باشند.
\begin{enumerate}
	\item
	نمونه‌‌های داده در فضای ورودی تصویر
	$\mathcal{X}=\{x^i\}$
	\item
	برچسب هر نمونه در سطح درشت‌دانه (
	\trans{ابردسته}{Superclass}
	)
	$\mathcal{Y}_{super}=\{y^i_{super}\}$
	\item 
	برچسب هر نمونه در سطح ریزدانه (
	\trans{دسته پایه}{Base Class}
	)
	$\mathcal{Y}_{base}=\{y^i_{base}\}$
	
\end{enumerate}

در سطح ابردسته فرض می‌کنیم 
$\mathcal{Y}_{super,target} \subseteq \mathcal{Y}_{super, source}$
که نشان می‌دهد برچسب‌های مقصد زیرمجموعه‌ی برچسب‌های دادگان مبدا است. همانطور که در 
\ref{sec:settingDef}
بیان شد، با توجه به وجود تغییر زیرجمعیت در مسئله‌ی دسته‌بندی بزرگ مقیاس، این موضوع رایج است که دسته‌های ریزدانه متفاوت باشند. اما در سطح درشت‌دانه با توجه به وجود سلسله‌مراتب ذاتی در داده‌ها، قادر به تعریف یک‌سری دسته‌ی کلی هستیم که بین دادگان مبدا و مقصد همواره مشترک‌اند. به عنوان مثال در گونه‌‌های حیوانات، ممکن است گونه‌ای به تازگی شناسایی شود اما دسته‌هایی در سطح پرندگان و آبزیان و ... همواره ثابت هستند. با توجه به این موارد، برای سنجش خاصیت مجموعه باز بودن و تحمل 
\trans{تغییر زیرجمعیت}{Sub-Population Shift}
، برچسب‌های سطح‌ریزدانه بین مجموعه دادگان مبدا و مقصد را ناسازگار و کاملا مجزا از هم در نظر گرفتیم و چالش‌برانگیزترین تنظیمات ممکن را اعمال کردیم 
$\mathcal{X}_{source}\cap \mathcal{X}_{target}=\emptyset$ و
$\mathcal{Y}_{base,source}\cap \mathcal{Y}_{base,target}=\emptyset$
.
یعنی هیچ اشتراکی بین نمونه‌‌ها و همچنین برچسب‌های ریزدانه وجود ندارد.

با توجه به اینکه قید اصلی مسئله، نمونه محدود بودن آن است، هر 
$y_{c}$
یعنی برچسب هر دسته، دارای تعدادی نمونه‌ی محدود چه در سمت دادگان مبدا و چه مقصد است. علاوه بر این، روش پیشنهادی نشان داده‌است که هیچ نیازی به داده‌ی آموزش زیاد نداریم چرا که با استفاده از مدل‌های بنیادی بزرگ توانسته‌ایم به 
\trans{داده‌کارایی}{Data Efficiency}
و داشتن ویژگی‌های بیانگر برسیم. با استفاده از این رویکرد، نشان داده‌ایم حتی در تنظیمات 
\trans{بدون آموزش}{Free Training}
توانسته‌ایم به کارایی مناسب دست یابیم.

برای یادگیری نمونه محدود، شبکه‌های پروتوتیپیکال
\cite{protonet}
را از بین رویکردهای متایادگیری انتخاب کرده‌ایم. این رویکرد در 
\ref{sec:protonet}
به طور کامل شرح داده‌شد. این انتخاب به دلیل مولد بودن این شبکه است و می‌تواند با دسته‌های پایه کاملا 
\trans{بدیع}{Novel}
در دادگان مقصد کار کند. فاز‌های متاآموزش و متاتست، به ترتیب روی دادگان مبدا و مقصد به صورت اپیزودیک انجام می‌شود. در هر اپیزود، یک مسئله‌ی یادگیری مستقل داریم که به عنوان یک وظیفه‌ی 
$\mathcal{T}$
از بین همه‌ی وظایف قابل تعریف روی 
$\mathcal{D}$
 نمونه برداشته ایم و در نتیجه در هر اپیزود، یک مسئله‌ی K 
 \trans{راه}{Way}
 N 
 \trans{شات}{Shot} 
 داریم. این نمادگذاری به معنی این است که در یک اپیزود K دسته داریم و به ازای هر دسته N نمونه‌ی پشتیبان وجود دارد که N تعداد محدودی نمونه است.
 
 در هر وظیفه، داده‌های مربوط به هر دسته به داده‌های پشتیبان و پرسش تقسیم ‌می‌شوند که به ازای یک دسته‌ی c آن‌ها را به ترتیب با 
 $\mathcal{S}^c$ 
 و
 $\mathcal{Q}^c$
 نمایش می‌دهیم. هدف متایادگیری رسیدن به دقت قابل قبول روی 
 $\mathcal{Q}$ 
 است که این هدف را با تطبیق سریع با به کار بردن نمونه‌های محدود پشتیبان 
  $\mathcal{S}$
  دنبال می‌کند.
  
  جدول نماد ها
  
  \begin{table}[htp]
  	\caption{هلو}
  	\centering
%  	\resizebox{\columnwidth}{!}{%
\begin{tabular}{@{}cc@{}}
	\toprule
	نماد                   & توضیح              \\ \midrule
	$\mathcal{D}_{source}$ & مجموعه دادگان مبدا \\
	$\mathcal{D}_{target}$ & مجموعه دادگان مقصد \\
	$\mathcal{X}$          &                    \\
	$\mathcal{Y}_{base}$   &                    \\
	$\mathcal{Y}_{super}$  &                    \\
	$\phi, \varphi$        &                    \\
	$\theta, \vartheta$    &                    \\
	$g$                    &                    \\
	$f$                    &                    \\
	$k-way$                &                    \\
	$N-shot$               &                    \\
	$\mathcal{T}$          &                    \\
	$\mathcal{S}^c$        &                    \\
	$\mathcal{Q}^c$        &                    \\ \bottomrule
\end{tabular}%
%  	}
  \end{table}

  
  ما از بدنه‌های بنیادی بزرگ برای تولید بازنمایی در مسئله‌ی دسته‌بندی نمونه محدود استفاده کرده‌ایم. این بدنه‌های بنیادی که وظیفه‌ی تولید تعبیه با وزن‌های 
  \trans{منجمد}{Freeze}
  را دارند، در سطح درشت‌دانه و ابردسته‌ها با 
  $g_\phi$
  و در سطح‌ریزدانه و دسته‌های پایه با 
  $g_\theta$
  نمایش داده شده‌اند. همچنین شبکه‌های کاملا متصلی که برای انجام دسته‌بندی بر روی بدنه‌ها سوار می‌شوند در سطح ابردسته‌ها و دسته‌های پایه به ترتیب با 
  $f_\varphi$
  و 
  $f_\vartheta$
  نشان داده‌ شده‌است که در واقع بخشی است که ما آموزش می‌دهیم. 
  
  به طور خلاصه، با توجه به بایسته‌ها و چالش‌های مسئله‌ی دسته‌بندی بزرگ مقیاس نمونه محدود، یک تنظیمات جدید در سخت‌ترین حالت ممکن تعریف کرده ایم که در سطح ریزدانه بین دادگان مبدا و مقصد هیچ اشتراکی از لحاظ داده و برچسب نداریم. از طرفی در سمت مبدا هم دادگان نمونه محدود است و مثل سایر روش‌ها نیازی به جمع‌آوری دادگان بزرگ در این مرحله نداریم. در فاز متاآزمون و داده‌های مقصد نیاز به برچسب‌گذاری در سطح درشت‌دانه را با اعمال یک تکنیک خاص از بین برده‌ایم و بنابراین روش‌ما نیازی به دانستن 
$\mathcal{Y}_{super, source}$
 برای داده‌های پشتیبان مقصد ندارد.

\section{جمع‌بندی روش و ارائه‌ی الگوریتم}


